\documentclass{article}
\usepackage{graphicx} % Required for inserting images
\usepackage{amsmath}
\usepackage{amsthm}
\usepackage{amsfonts}
\usepackage{cleveref}
\usepackage{aliascnt}
\usepackage{dsfont}

\usepackage[
    backend=biber,
    style=authoryear,
    sorting=nyt,
    url=false,
    dashed=false,
    giveninits=true,
    maxnames=99,
    maxcitenames=2,
    mincitenames=1
]{biblatex}

\DeclareNameAlias{author}{last-first}

\addbibresource{ref.bib}

\newtheorem{definition}{Definition}[section]

\newaliascnt{theorem}{definition}
\newtheorem{theorem}[theorem]{Theorem}
\aliascntresetthe{theorem}

\crefname{definition}{definition}{definitions}
\Crefname{definition}{Definition}{Definitions}

\crefname{theorem}{theorem}{theorems}
\Crefname{theorem}{Theorem}{Theorems}

\title{MEGPD Markov Process: properties and their proofs}

\author{Andrea Ferrero}

\begin{document}

\maketitle

\section{Model definition}

The Markov process is defined on a continuous state space by the transition density $p(y|x) = p(x,y) / p_X(x)$ where the joint density is that of the multivariate EGPD (MEGPD) and the common marginal $p_X(x)$ is the implied Extended GPD (EGPD).

More precisely, the EGPD distribution function is defined by:

\begin{equation*}
  F(x) = B(H_\xi(x)^\kappa),
\end{equation*}

where $H_\xi$ is the GPD, $\kappa$ is a positive parameter and $B$ is an absolutely continuous cdf on $[0, 1]$ whose derivative satisfies: $0<b(0)<\infty$ and $0 < b(1) < \infty$.

The density $f(x)$ of an EGPD is:

\begin{equation*}
  f(x) = \kappa h_\xi(x) H_\xi(x)^{\kappa-1} b(H_\xi(x)^\kappa).
\end{equation*}

Recall that the EGPD is a regularly varing random variable both at infinity and at zero, more precisely $X \in RV_{\frac{1}{\xi}}^\infty$ and $X \in RV_\kappa^0$. Regular variation at infinity and at zero come from the conditions imposed on the $B$ function.

Namely, we have that:

\begin{equation*}
  \lim_{x \to \infty} \frac{f(x)}{\kappa h_\xi(x)} = b(1),
\end{equation*}

and for the lower tail:

\begin{equation*}
  \lim_{x \to 0} \frac{f(x)}{\kappa x^{\kappa-1}} = b(0).
\end{equation*}

From the density behaviours it follows that the density of an EGPD is $RV_{-\frac{1}{\xi} - 1}^\infty$ and $RV_{\kappa - 1}^0$.

The join density of the MEGPD is then defined on the radial/angular decomposition of the vector, leading to:

\begin{equation*}
  p(x, y) = p_R(x + y) \phi_{\delta(x+y)}(\log \frac{x}{y}) \frac{x+y}{xy},
\end{equation*}

where $p_R$ is another EGPD density (note that the sum of EGPD leads to an EGPD with the same tail index $\xi$) while $\phi_\delta$ is the density of the zero mean normal with variance $\delta(x+y)^2$. By assumption, we have that $\delta(x+y) \to \delta_0 > 0$, $\delta(x+y) \to \delta_\infty$ for the radius $r = x + y$ going to zero and infinity, respectively.

Note that both the joint and the marginal densities are continuous and positive on $(0, \infty)^2$ and $(0, \infty)$, respectively.

\section{Ergodicity of the MEGPD Markov Process}

To fix notation, we generally work on a measurable space $(E, \mathcal{E})$, where in our case $E = (0, \infty)$ and $\mathcal{E}$ is the Borel $\sigma$-algebra on $E$.

\begin{definition}[Markov transition kernel]
  A Markov transition kernel is a function
  \[
    \Pi : E \times \mathcal{E} \to [0,1]
  \]
  such that:
  \begin{enumerate}
    \item For every $x \in E$, the function $A \mapsto \Pi(x,A)$ is a probability measure on $(E,\mathcal{E})$.
    \item For every $A \in \mathcal{E}$, the function $x \mapsto \Pi(x,A)$ is measurable.
  \end{enumerate}
  The quantity $\Pi(x,A)$ represents the probability of moving to a measurable set $A$ from a starting state $x$.
\end{definition}

In our case, the transition probability admits a density $p(y \mid x)$ with respect to the Lebesgue measure, so that

\[
  \Pi(x,A) = \int_A p(y \mid x)\,dy.
\]

We denote by $\Pi^n(x,A)$ the probability of starting from $x$ and landing in $A$ after $n$ steps.

Using the MEGPD to construct a Markov process of order one, we can show that the resulting process is ergodic.

Recall that, in order for a Markov process to be ergodic, we need to show that it admits a stationary distribution and that the process is irreducible and aperiodic.

We refer to \textcite[Theorem 4]{roberts2004} where the only assumption is that the $\sigma$-algebra needs to be countably generated, which is the case for the usual Borel $\sigma$-algebra we are considering as mentioned by the authors. The theorem is stated under $\pi-a.s.$ (with $\pi$ the stationary distribution) but later it is mentioned that for absolutely continuous transition distributions with respect to $\pi$ the theorem holds for all $x \in E$. More generally, Harris recurrence is invoked.

\subsection{Stationarity}

We can show that the Markov process defined by the MEGPD admits a stationary distribution $\pi$, which coincides with the common marginal EGPD $p_X$.

The key observation is that the joint distribution is symmetric. Indeed, it is constructed from $p_R(x+y)$, which is symmetric in $x$ and $y$, the normal density $\phi_{\delta(x+y)}(\log(x/y))$, which is symmetric by definition, and the Jacobian product, which is symmetric as well.

It therefore suffices to establish the detailed balance condition.

\begin{definition}[Stationary distribution]
  A probability measure $\pi$ on $(E,\mathcal{E})$ is said to be stationary for the Markov transition kernel $\Pi$ if
  \[
    \pi(A) = \int_E \Pi(x,A)\,\pi(dx),
    \qquad A \in \mathcal{E}.
  \]
  If $\pi$ admits a density with respect to the Lebesgue measure, we denote this density by $\pi(x)$.

\end{definition}

\begin{theorem}[Detailed balance and stationarity]
  Let $\Pi$ be a Markov transition kernel on $(E,\mathcal{E})$, admitting a transition density $p(y \mid x)$ with respect to the Lebesgue measure. If a probability density $\pi$ satisfies the detailed balance condition
  \[
    \pi(x)p(y \mid x)
    =
    \pi(y)p(x \mid y)
  \]
  for almost every $(x,y) \in E^2$, then $\pi$ is a stationary density for the Markov process.
\end{theorem}

\begin{proof}
  By definition, a density $\pi$ is stationary if
  \[
    \int_E \pi(x)p(y \mid x)\,dx = \pi(y).
  \]
  Using the detailed balance condition, we obtain
  \[
    \begin{aligned}
      \int_E \pi(x)p(y \mid x)\,dx
       & = \int_E \pi(y)p(x \mid y)\,dx \\
       & = \pi(y)\int_E p(x \mid y)\,dx \\
       & = \pi(y),
    \end{aligned}
  \]
  where the last equality follows from the fact that $p(\cdot \mid y)$ is a probability density. Therefore, $\pi$ is stationary.
\end{proof}

\begin{theorem}[Stationarity of the MEGPD Markov process]
  Suppose that the joint density $p(x,y)$ of the MEGPD is symmetric, that is,
  \[
    p(x,y) = p(y,x).
  \]
  Then the common marginal density $p_X$ is a stationary distribution for the Markov process.
\end{theorem}

\begin{proof}
  Since the joint distribution admits the factorization
  \[
    p(x,y) = p_X(x)p(y \mid x),
  \]
  its symmetry gives
  \[
    \begin{aligned}
      p_X(x)p(y \mid x)
       & = p(x,y)             \\
       & = p(y,x)             \\
       & = p_X(y)p(x \mid y).
    \end{aligned}
  \]
  Hence, the detailed balance condition holds with $\pi = p_X$ and $p_X$ is stationary.
\end{proof}

\subsection{Irreducibility}

We now show that the Markov process is irreducible with respect to the Lebesgue measure $\lambda$.

\begin{definition}[$\lambda$-irreducibility]
  A Markov process with transition kernel $\Pi$ is said to be $\lambda$-irreducible if, for every $x \in E$ and every measurable set $A \in \mathcal{E}$ satisfying $\lambda(A)>0$, there exists an integer $n \geq 1$ such that
  \[
    \Pi^n(x,A)>0.
  \]
\end{definition}

In our case, irreducibility follows directly from the assumption that both the joint and marginal densities of the Markov process are positive on $E=(0,\infty)$. Consequently, the conditional density $p(y \mid x)$ is also positive on $E$.

\begin{theorem}[Irreducibility of the MEGPD Markov process]
  Suppose that the conditional density satisfies
  \[
    p(y \mid x)>0,
    \qquad x,y \in E.
  \]
  Then the Markov process is $\lambda$-irreducible.
\end{theorem}

\begin{proof}
  It is sufficient to consider $n=1$. Let $x \in E$ be arbitrary and let $A \in \mathcal{E}$ satisfy $\lambda(A)>0$.

  For every $m \geq 1$, define
  \[
    A_m =
    \left\{
    y \in A :
    p(y \mid x) \geq \frac{1}{m}
    \right\}.
  \]
  Since the conditional density is strictly positive, we have
  \[
    A = \bigcup_{m=1}^{\infty} A_m.
  \]
  As $\lambda(A)>0$, there must exist some $m \geq 1$ such that $\lambda(A_m)>0$.

  Therefore,
  \[
    \begin{aligned}
      \Pi(x,A)
       & = \int_A p(y \mid x)\,dy        \\
       & \geq \int_{A_m} p(y \mid x)\,dy \\
       & \geq \frac{1}{m}\lambda(A_m)    \\
       & > 0.
    \end{aligned}
  \]
  Hence, the Markov process is $\lambda$-irreducible.
\end{proof}

Irreducibility also guarantees that the stationary distribution is unique, provided the process admits a stationary probability distribution.

\subsection{Small Sets}

For aperiodicity, as well as for later results, we need to introduce the notion of a small set.

\begin{definition}[$n$-small set]
  Let $\Pi$ be a Markov transition kernel on $(E,\mathcal{E})$. A measurable set $C \in \mathcal{E}$ is said to be $n$-small if there exist a nonzero finite measure $\nu$ on $(E,\mathcal{E})$ and a constant $\varepsilon>0$ such that
  \[
    \Pi^n(x,A)
    \geq \varepsilon \nu(A),
    \qquad x \in C,\ A \in \mathcal{E}.
  \]
\end{definition}

Now we want to prove that compact sets are small sets if the conditional density is positive.

\begin{theorem}[Compact sets are small]
  \label{thm:compact-small}
  Suppose that the conditional density $p(y \mid x)$ is strictly positive and continuous on $E \times E$, where $E=(0,\infty)$.

  Then, for any compact sets $C,K \subset E$ with $\lambda(K)>0$, the set $C$ is $1$-small.
\end{theorem}

\begin{proof}
  Let $C \subset E$ be compact and choose a compact set $K \subset E$ such that $\lambda(K)>0$.

  Since $C$ and $K$ are compact, their Cartesian product $C \times K$ is also compact. Moreover, by assumption, the conditional density $p(y \mid x)$ is continuous and strictly positive on $E \times E$.

  Therefore, by the Extreme Value Theorem, there exist $(x_0, y_0) \in C \times K$ such that

  \[
    m
    =
    \min_{(x,y)\in C\times K} p(y \mid x) = p(y_0 \mid x_0) > 0,
  \]
  where we used the fact that the conditional density is positive.

  Consequently, for every $x\in C$ and every $A\in\mathcal{E}$,
  \[
    \begin{aligned}
      \Pi(x,A)
       & = \int_A p(y \mid x)\,dy           \\
       & \geq \int_{A\cap K}p(y \mid x)\,dy \\
       & \geq m \int_{A\cap K}dy            \\
       & = m \lambda(A\cap K).
    \end{aligned}
  \]

  Defining the measure
  \[
    \nu(A)=\frac{\lambda(A\cap K)}{\lambda(K)},
    \qquad A\in\mathcal{E},
  \]
  we obtain a probability measure supported on $K$. Hence,
  \[
    \Pi(x,A)
    \geq m \lambda(K)\nu(A),
    \qquad x\in C,\ A\in\mathcal{E}.
  \]

  Thus, taking $\varepsilon = m \lambda(K)>0$, we obtain
  \[
    \Pi(x,A) \geq \varepsilon \nu(A),
    \qquad x \in C,\ A \in \mathcal{E}.
  \]
  Therefore, $C$ is $1$-small.
\end{proof}

\subsection{Aperiodicity}

From \textcite[Definition D.2.3]{kulik2020}, we can see a convenient definition of aperiodicity in terms of small sets:

\begin{definition}[Aperiodicity]
  A markov kernel $\Pi$ which admits an accessible small set $C$, i.e. a set for which the return times are almost surely finite, is said to be aperiodic if there exists no integer $c > 1$ such that:

  \begin{equation*}
    \{n \in \mathbb{N}^* : inf_{x \in C} \Pi^n (x, C) > 0\} \subset \{c, 2c, 3c, \dots\} .
  \end{equation*}
\end{definition}

Moreover, they mention a convenient criterion for aperiodicity: the kernel $\Pi$ is aperiodic if and only if there exists an \textbf{increasing countable} union of small sets. \parencite[Lemma D.2.4]{kulik2020}

Incidentally, if a Markov process has an accessible small set it is said to be irreducible (proving in a different way that the MEGPD Markov process is irreducible).

\begin{theorem}[Aperiodicity of the MEGPD Markov process]
  Suppose that the conditional density $p(y\mid x)$ is strictly positive and continuous on $E\times E$, where $E=(0,\infty)$.

  Then the Markov kernel $\Pi$ is aperiodic.
\end{theorem}

\begin{proof}
  From \cref{thm:compact-small}, every compact subset of $E$ is small.

  Consider the sequence of compact intervals
  \[
    C_n=\left[\frac{1}{n},n\right],
    \qquad n\geq 1.
  \]

  Clearly, each $C_n$ is compact and the sequence is increasing, since
  \[
    C_n\subseteq C_{n+1},
    \qquad n\geq 1.
  \]

  Furthermore, we have
  \[
    \bigcup_{n=1}^{\infty}C_n
    =(0,\infty)=E.
  \]

  Since every $C_n$ is compact, it is also small by the preceding result. We have therefore constructed an increasing countable union of small sets covering the entire state space.

  Hence, by \textcite{kulik2020}, Lemma D.2.4, the Markov kernel $\Pi$ is aperiodic.
\end{proof}

\subsection{Ergodicity}

\begin{definition}[Total variation distance]
  The total variation distance between two probability measures $\nu_1(\cdot)$ and $\nu_2(\cdot)$ is:

  \[
    \left\| \nu_1(\cdot) - \nu_2(\cdot) \right\|_{\mathrm{TV}} = \sup_A \mid \nu_1(A) - \nu_2(A) \mid
  \]
\end{definition}

\begin{theorem}[Ergodicity of the MEGPD Markov process]
  \label{thm:megpd-ergodicity}
  Suppose that the joint density of the MEGPD model is strictly positive and continuous on $E\times E$, where $E=(0,\infty)$, and symmetric with respect to its arguments.

  Then the Markov process constructed from the MEGPD model is irreducible and aperiodic, with unique stationary distribution given by the common marginal EGPD density $p_X$.

  Moreover, for any initial state $X_0=x\in E$,
  \[
    \lim_{n\to\infty}
    \left\|\Pi^n(x,\cdot)-\pi(\cdot)\right\|_{\mathrm{TV}}
    =0,
  \]
  where $\pi(dx)=p_X(x)\,dx$.

  In particular, $\lim_{n \to \infty} \Pi^n(x, A) = \pi(A)$ for all measurable $A \subseteq \mathcal{E}$
\end{theorem}

\section{Geometric Ergodicity}
After having shown that the Markov process is ergodic, we might want to know how long does it take for the chain to reach the stationary distribution.

\begin{definition}[Geometric Ergodicity]
  A Markov process with stationary distribution $\pi$ is geometrically ergodic if:

  \[
  \left\| \Pi^n(x, \cdot) - \pi(\cdot) \right\| \leq M(x) \rho^n
  \]

  for some $\rho < 1$, where $M(x) < \infty$ for $\pi-a.e.$ $x \in E$. 
\end{definition}

To ensure geometric ergodicity, we need to satisfy a certain drift condition:

\begin{definition}[Drift condition]
  The Markov process is said to satisfy a drift condition for a given set $C$ if there are constants $\lambda (0, 1)$ and $b < \infty$ and a function $V: E \to [1, \infty]$ such that:
  \[
    PV \leq \lambda V + b \mathds{1}_C ,
  \]
  where the function $PV$ refers to the conditional expectation $\mathbb{E}\left[V(y) \mid X = x\right]$
\end{definition}

\begin{theorem}[Drift condition, small sets and geometric ergodicity]
  Conider a $\phi$-irreducible, aperiodic Markov process with stationary distribution $\pi$. Suppose that for a small set $C \subset E$ the drift condition holds. Then the process is geometrically ergodic.
\end{theorem}

We are going to try to prove the drift condition for the MEGPD Markov process focusing on compact sets thanks to \cref{thm:compact-small}.

\subsection{Current proof status of Geometric Ergodicity}

We need to show that $\mathbb{E}[V(y) \mid X = x] \leq \lambda V(x) + \mathds{1}_C$ with $V \geq 1$, $\lambda \in (0, 1)$ for some choice of $V$.

We can 

\end{document}